\subsection{可用容量消融的账单与实施核验}\label{app:capacity}
三档窗口保持初末6,000 kWh与5,000 kW功率一致；对0和4,800 kWh窗口的四机制均从2月1日起重新执行完整334日。9,600 kWh窗口使用与主表、工作簿相同的固定终点轨迹，并用同一入口从保存的12月30日库存重算真实末端两日；四机制的原合同、最终合同、充放电、紧急电、弃电、库存及结算电价逐项与主轨迹一致。

计算中，窗口边界同时进入合同LP、库存价值递推与可达执行器；每个窗口在独立进程中运行。零窗口的物理动作由停用分支固定为零，单点状态无需构造差分斜率。半窗口与全窗口保持30 kWh节点间距，其余主参数不重新选择。每次预报场景的训练目标截止时刻不晚于规划时刻；新增窗口沿用主策略的信息权限。
\begin{table}[H]\centering\small
\caption{三档可用容量的完整运行期账单}\label{tab:capacitycost}
\begin{tabular}{lrrrrrr}\toprule
机制& $W=0$ & $W=4.8$ & $W=9.6$ & 全对零节省 & 降幅 & 全对半节省\\\midrule
问题二&1,881.03&1,557.14&1,384.42&496.61&26.40\%&172.72\\
问题三&1,864.20&1,522.42&1,346.18&518.01&27.79\%&176.23\\
问题四日前&1,952.47&1,638.56&1,465.46&487.01&24.94\%&173.10\\
问题四日内&1,934.98&1,601.66&1,425.21&509.77&26.35\%&176.45\\
\bottomrule\end{tabular}
\par\footnotesize $W$单位MWh，费用单位万元。全对零节省为$C(0)-C(9.6)$，降幅除以$C(0)$；全对半节省为$C(4.8)-C(9.6)$。
\end{table}

\begin{table}[H]\centering\small
\caption{容量消融费用分解（万元）}\label{tab:capacityparts}
\begin{tabular}{lrrrrrr}\toprule
机制& $W$/MWh & 原计划 & 增购 & 退购净额 & 紧急费 & 合计\\\midrule
问题二&0.0&1,749.58&0.00&0.00&131.45&1,881.03\\
问题二&4.8&1,488.46&0.00&0.00&68.68&1,557.14\\
问题二&9.6&1,319.71&0.00&0.00&64.70&1,384.42\\
问题三&0.0&1,759.12&25.28&-20.17&99.98&1,864.20\\
问题三&4.8&1,502.91&24.58&-28.12&23.05&1,522.42\\
问题三&9.6&1,334.07&25.22&-30.57&17.47&1,346.18\\
问题四日前&0.0&1,813.57&0.00&0.00&138.90&1,952.47\\
问题四日前&4.8&1,567.20&0.00&0.00&71.36&1,638.56\\
问题四日前&9.6&1,399.98&0.00&0.00&65.48&1,465.46\\
问题四日内&0.0&1,823.77&32.95&-21.28&99.53&1,934.98\\
问题四日内&4.8&1,580.41&28.43&-30.49&23.30&1,601.66\\
问题四日内&9.6&1,411.77&28.73&-32.77&17.48&1,425.21\\
\bottomrule\end{tabular}
\par\footnotesize 分项按未舍入数值相加；退购净额为退款扣除退购手续费后的负支出。
\end{table}

日差按较窄窗口费用减较宽窗口费用计算。下表使用全部重叠、非环绕7日块，有放回抽取并拼接至334日，重复10,000次。种子序列为$[20260912,k,p]$：机制索引$k=0,1,2,3$，停用至全窗口比较的$p=1$，半窗口至全窗口的$p=2$。每日正负差与月度累计均保存在对应逐日账单和核验文件中；不将重采样重复解释为独立运行年份。
\begin{table}[H]\centering\small
\caption{容量配对日差的7日块重采样（元/日）}\label{tab:capacityci}
\begin{tabular}{llrrr}\toprule
机制&比较&日均节省&95\%区间&节省/增支/持平日\\\midrule
问题二&停用至全窗口&14,868.51&[14,464.03, 15,266.46]&334/0/0\\
问题二&半窗口至全窗口&5,171.37&[5,062.54, 5,281.57]&334/0/0\\
问题三&停用至全窗口&15,509.32&[15,296.67, 15,775.95]&334/0/0\\
问题三&半窗口至全窗口&5,276.49&[5,197.99, 5,360.38]&334/0/0\\
问题四日前&停用至全窗口&14,581.08&[14,092.41, 15,089.54]&333/1/0\\
问题四日前&半窗口至全窗口&5,182.56&[5,027.95, 5,330.61]&334/0/0\\
问题四日内&停用至全窗口&15,262.66&[14,859.61, 15,716.02]&334/0/0\\
问题四日内&半窗口至全窗口&5,282.91&[5,124.99, 5,433.48]&334/0/0\\
\bottomrule\end{tabular}
\par\footnotesize 正差为较大窗口降低费用。块长7日，10,000次，种子序列$[20260912,k,p]$；当前单年度时间稳定性诊断。
\end{table}

12组轨迹的库存区间、初末条件、跨午夜衔接、充放电功率、单向动作、紧急电量和供需平衡均通过独立逐段检查。电量项容差为$10^{-5}$ kWh，功率按每段$5000/6$ kWh约束核验。现金费用同时按分项账单和逐段$p_t\{r_t+0.5|r_t-q_t|+5e_t\}$重算，年度总费与保存结果的绝对差小于$10^{-5}$元。实际价格逐段与题目数据对应，日前机制无日内改约，日内机制的每次发布只覆盖当前及后续区间。所有新容量运行的源码、输入和实验入口哈希均与审计观察快照相符；主结果五份工作簿保持原对应轨迹。


\paragraph{零窗口日前合同的解析复核}
在每日0时，两个日前机制满足$r_t=q_t$，当前日合同变量彼此独立；零库存增量使状态不再连接各段购电。当段有限场景目标为
\begin{equation}\label{eq:zeroquantile}
 f_t(q)=\frac1S\sum_{s=1}^S p_t^s\{q+5(n_t^s-q)_+\},\qquad q\ge0.
\end{equation}
在正场景电价下，令$\omega_s=p_t^s/\sum_i p_t^i$，对应净需求的加权分布函数为$F_\omega$。$f_t$的左右导数分别为$(\sum_s p_t^s-5\sum_{n_t^s\ge q}p_t^s)/S$和$(\sum_s p_t^s-5\sum_{n_t^s>q}p_t^s)/S$，故无非负边界时最优条件为$F_\omega(q^-)\le0.8\le F_\omega(q)$。取一个加权0.8分位数并截断至非负，即得到可选最优合同。当前日的合同系数与未来日增益分组独立，零窗口的库存等式使未来段费用不再影响当段合同，因此可逐段检验。分位数最优解可能不唯一，核验以同场景目标差为准。
针对两个日前机制各48,096段，在每日0时的相同场景下，保存合同与解析分位合同的单段目标最大绝对差分别约为$3.62\times10^{-10}$元和$3.59\times10^{-10}$元。该核验验证零窗口下当前有限场景规划的实现，比较对象是规划目标，而非未知真实分布下的最优费用。

